\section{Background on variational inference for latent SDEs}\label{app:sec:background}

In this section, we provide background on variational inference for latent SDEs. 
First, in \Cref{app:subsec:assumptions} we describe our formal assumptions necessary for the probabilistic model and variational inference objective to be well-posed.   
In \Cref{app:subsec:elbo} we derive the evidence lower bound (ELBO) on the observed data log likelihood.
In \Cref{app:subsec:div-free} we provide further discussion of divergence-free vector fields and prove \Cref{prop:div_free}.
In \Cref{app:subsec:gaussian} we characterize the affine divergence-free vector fields for Gaussian one-time marginals (\Cref{app:prop:gaussian-compatible-drift}).
We conclude in \Cref{app:subsec:svise-sde-compare} by discussing two simulation-free VI algorithms, SVISE (\citet{course2023state, course2023amortized}) and SDE Matching (\citet{bartosh2025sde}). For each algorithm we derive the choice of posterior drift and discuss how this choice relates to Helmholtz-SDE.

For full generality, we allow the prior drift $\mbf_p$ to be time-dependent and for the diffusion coefficient $\mbSigma$ to be both space- and time-dependent.

\subsection{Technical Assumptions}\label{app:subsec:assumptions}
 
In this subsection, we detail our formal assumptions on the prior \eqref{eqn:prior_sde} and posterior \eqref{eqn:variational_family} processes. 
The assumptions are stated abstractly here and are verified for the variational family considered by Helmholtz-SDE in \Cref{app:subsec:verify-assumptions}.
They are required for every admissible choice of model parameters $\mbtheta$ and variational parameters $\mbphi$.  
Here, $\mbf_q$ denotes the \emph{full} posterior drift in \cref{eqn:variational_family}; for Helmholtz-SDE, this is the sum of the reference drift and the Helmholtz correction.

\paragraph{Existence and uniqueness.}
We first require that the prior and posterior SDEs, \cref{eqn:prior_sde} and \cref{eqn:variational_family}, respectively, admit unique strong solutions.
 
\begin{assumption}[SDE existence and uniqueness]\label{app:assum:regularity}
The prior and posterior SDEs admit unique strong solutions on $[0,T]$.
\end{assumption}

For the prior SDE, we follow the classical route of imposing Lipschitz continuity and linear growth on the drift and diffusion coefficient (see, e.g., \citet[Theorem 5.2.1]{oksendal2013stochastic}): there exist constants $C_1, C_2 > 0$ such that for all $\mbx, \mby \in \bbR^K$ and $t \in [0, T]$
\begin{equation*}
\begin{aligned}
    &\|\mbf_p(\mbx,t)- \mbf_p(\mby,t)\|_2 + \| \mbSigma(\mbx, t)^{1/2} -  \mbSigma(\mby, t)^{1/2} \|_F \leq  C_1\|\mbx - \mby\|_2, \quad &\text{(Lipschitz)}\\ 
    &\|\mbf_p(\mbx,t)\|_2 + \|\mbSigma(\mbx, t)^{1/2}\|_F \leq  C_2(1+\|\mbx\|_2). \quad &\text{(linear growth)}
\end{aligned}
\end{equation*}
For example, if $\mbf_p$ is a neural network that takes as input $\mbx$ and $t$, then Lipschitz continuity and linear growth can be ensured by restricting the operator norm of the weight matrices and choosing (globally) Lipschitz activation functions. 

For the posterior SDE, the polynomial approximation to the Helmholtz correction will not satisfy linear growth when the polynomial degree is $\ell > 1$.
We will therefore pursue an alternate route to ensure existence and uniqueness.

\paragraph{Densities.}
We additionally require that, for every $0 \leq t \leq T$, the one-time marginals of $\bbP$ and of any $\bbQ \in \cQ$ are absolutely continuous with respect to the $K$-dimensional Lebesgue measure. The corresponding densities are written as $\{p(\mbx, t)\}_{0 \leq t \leq T}$ and $\{q(\mbx, t)\}_{0 \leq t \leq T}$, respectively. 
 
\begin{assumption}[Marginal densities]\label{app:assum:densities}
For every $t\in[0,T]$, the one-time marginals of $\bbP$ and $\bbQ$ admit Lebesgue densities.
Moreover, $\mbSigma(\mbx,t)$ is positive-definite for all $(\mbx,t)\in\bbR^K\times[0,T]$.
\end{assumption}
 
\paragraph{Absolute continuity.}
In order for the upper bound to the negative marginal log likelihood (negative ELBO) to be well-defined, we need the approximate posterior SDE $\bbQ \in \cQ$ to possess a probability density with respect to the prior $\bbP$ and for the KL divergence between $\bbQ$ and $\bbP$ to be finite. 
These two requirements are the subjects of the following two assumptions.
 
\begin{assumption}[Absolute continuity]\label{app:ass:path-ac}
The posterior SDE \cref{eqn:variational_family} admits a probability density $d\bbQ/d\bbP(\mbx)$ with respect to the prior SDE.
\end{assumption}

\begin{assumption}[Finite relative entropy]\label{app:ass:finite-kl}
   The initial distributions satisfy
   \begin{equation*}
       \KL{q(\mbx,0)}{p(\mbx,0)}<\infty
   \end{equation*}
    and the difference between the prior and posterior drifts satisfies
    \begin{equation*}
        \int_0^{T} \bbE_{\bbQ}\| \mbSigma(\mbx(t), t)^{-1/2}\{\mbf_p(\mbx(t),t) -  \mbf_q(\mbx(t), t)\}\|_2^2 dt < \infty.
    \end{equation*}
\end{assumption}
 

Under these assumptions, we can obtain an expression for the probability density of $\bbQ$ with respect to $\bbP$ via Girsanov's theorem (see e.g., \citet[Corollary 3.5.13]{karatzas2014brownian}).
\begin{theorem}[Girsanov's theorem for SDEs]\label{app:thm:girsanov}
Suppose \Cref{app:assum:regularity,app:assum:densities,app:ass:path-ac,app:ass:finite-kl} hold. Then the probability density of $\bbQ \in \cQ$ with respect to $\bbP$ is
{\small
\begin{equation*}
\begin{aligned}
&\frac{d \bbQ}{d \bbP}(\mbx) = \frac{ q(\mbx, 0)}{p(\mbx, 0)}\exp\left\{ \int_0^{T} \sum_{i=1}^{K} \left[\mbSigma^{-1/2} \mbb(\mbx(t), t)\right]_i d\mbw_i(t) - \frac{1}{2} \int_0^{T} \sum_{i=1}^{K} \left[\mbSigma^{-1/2} \mbb(\mbx(t), t)\right]_i^2 dt \right\},\\
&\mbb(\mbx, t) = \mbf_q(\mbx, t) - \mbf_p(\mbx, t),
\end{aligned}
\end{equation*}
}
where $(\mbw(t))_{0 \leq t \leq T}$ is a $K$-dimensional Brownian motion with respect to $\bbP$. 

\smallskip
In particular, for any bounded, measurable function $\Phi$ defined on $K$-dimensional continuous functions $\mbx$ having domain $[0, T]$,
\begin{equation*}
    \bbE_{\bbQ}[\Phi(\mbx)] = \bbE_{\bbP}\left[\Phi(\mbx) \frac{d \bbQ}{d \bbP}(\mbx) \right].
\end{equation*}
\end{theorem}

As a direct corollary of \Cref{app:thm:girsanov}, we obtain an expression for the KL divergence between the posterior and prior
{\small
\begin{align}
    &\KL{\bbQ}{\bbP} = \bbE_{\bbQ}\left[\log\left(\frac{d\bbQ}{d\bbP}(\mbx)\right)\right]\nonumber\\
    &= \KL{q(\mbx, 0)}{p(\mbx, 0)} + \frac{1}{2} \int_0^{T} \bbE_{\bbQ}\| \mbSigma(\mbx(t), t)^{-1/2}\{\mbf_p(\mbx(t),t) -  \mbf_q(\mbx(t), t)\}\|_2^2 dt. \label{app:eqn:path-kl}
\end{align}
}
The stochastic integral term vanishes because it is a mean-zero martingale.

\paragraph{Fokker-Planck equation.}
For a drift $\mbf$ and initial density $\rho(\mbx,0)$, we call $\{\rho(\mbx,t)\}_{0\leq t\leq T}$ a \emph{weak solution} to the Fokker--Planck equation if for every infinitely differentiable function $\varphi: \bbR^K \to \bbR$  having compact support,
\begin{equation*}
    t \longmapsto \int \varphi(\mbx)\rho(d\mbx, t)
\end{equation*}
is absolutely continuous and 
\begin{equation}\label{app:eqn:fpk-weak}
    \frac{d}{dt} \int \varphi(\mbx) \rho(d\mbx, t) = \int\left( \mbf(\mbx, t) \cdot \nabla_{\mbx} \varphi(\mbx) + \frac{1}{2} \mathrm{tr}(\mbSigma(\mbx, t) \nabla_{\mbx}^2 \varphi(\mbx)) \right) \rho(d\mbx, t), \quad \text{a.e.~$t$}.
    \end{equation}

Since $\bbQ \in \cQ$ is a solution to the SDE \cref{eqn:variational_family},  we know that $\{q(\mbx, t)\}_{0 \leq t \leq T}$ is a weak Fokker--Planck solution.
However, in order to argue that two variational SDEs whose drifts differ by a divergence-free vector field induce the same marginals, we need the solution to the Fokker--Planck equation to be unique within a suitable regularity class.
    
\begin{assumption}[Uniqueness of the Fokker--Planck solution]\label{app:ass:fpk-unique}
For every member $\bbQ \in \cQ$ of the variational family having drift $\mbf_q$, the marginal densities belong to the class
\begin{equation}\label{app:eqn:fpk-regularity-class}
        \int_0^T\int \left( \|\mbf_q(\mbx,t)\|_2 + \|\mbSigma(\mbx,t)\|_F \right) \rho(d\mbx,t) dt <\infty.
\end{equation}
The associated weak Fokker--Planck equation with initial law $q(\mbx, 0)$ has at most one solution belonging to the class \cref{app:eqn:fpk-regularity-class}.
\end{assumption}
    
\subsection{ELBO derivation}\label{app:subsec:elbo}
In this subsection, we derive the evidence lower bound (ELBO) in the latent SDE model. 

Adopting the notation from \Cref{app:thm:girsanov}, we have
\begin{equation*}
\begin{aligned}
    &\log p(\mby) = \log \int p(\mby | \mbx) \bbP(d\mbx)\\
    &= \log \int p(\mby | \mbx) \frac{d \bbP(\mbx)}{d \bbQ(\mbx)} \bbQ(d\mbx)\\
    &\overset{(i)}{\geq}  \int \log \left( p(\mby | \mbx)  \frac{d \bbP(\mbx)}{d \bbQ(\mbx)} \right) \bbQ(d\mbx)
    \\
    &= \bbE_{\bbQ}[\log p(\mby | \mbx)] - \KL{\bbQ}{\bbP}\\
    &\overset{(ii)}{=} \bbE_{\bbQ}[\log p(\mby | \mbx)] - \KL{q(\mbx, 0)}{p(\mbx, 0)} - \frac{1}{2} \int_0^{T} \bbE_{\bbQ}\| \mbSigma^{-1/2}\{\mbf_p(\mbx(t),t) -  \mbf_q(\mbx(t),t)\}\|_2^2 dt.
\end{aligned}
\end{equation*}
Equality (i) is an application of Jensen's inequality and equality (ii) is an application of Girsanov's theorem. Taking the negative of this expression yields the negative ELBO \cref{eqn:elbo}. 

Notice that each term in the negative ELBO depends both on the parameters of the generative model $\mbtheta$ and on the parameters of the variational posterior $\mbphi$.

\subsection{Divergence-free vector fields}\label{app:subsec:div-free}
In this subsection, we discuss the role of $q$-divergence-free vector fields in determining the expressivity of the variational family $\cQ$. We prove \Cref{prop:div_free}, which characterizes the non-uniqueness in path law that arises from marginal-based parameterizations of the variational posterior.

In \cref{eqn:fpk}, the Fokker-Planck equation is stated as a pointwise equality for $(\mbx, t) \in \bbR^K \times [0, T]$; this is only valid under regularity assumptions on the posterior SDE beyond those in \Cref{app:subsec:assumptions}.
Still, we can state the divergence-free condition in its \emph{weak form}, i.e., using the weak form of the Fokker--Planck equation \cref{app:eqn:fpk-weak}
\begin{equation}\label{app:eqn:div-free-weak}
    \int \mbv(\mbx, t) \cdot \nabla_{\mbx} \varphi(\mbx)  q( d\mbx, t)  = 0 \quad \text{a.e.~$t$}
\end{equation}
for all infinitely differentiable test functions $\varphi$ having compact support. Indeed, if $\mbv(\mbx, t)$ and $q(\mbx, t)$ are smooth then integration by parts implies
\begin{equation*}
    0 = \int \mbv(\mbx, t) \cdot \nabla_{\mbx} \varphi(\mbx)  q(d\mbx, t) = - \int \nabla_{\mbx} \cdot[ q( \mbx, t) \mbv(\mbx, t)] \varphi(\mbx)  d \mbx \quad \text{a.e.~$t$}
\end{equation*}
for all test functions $\varphi$. By Fubini's theorem, this implies $\nabla \cdot[ q( \mbx, t) \mbv(\mbx, t)] = 0$ a.e.~$(\mbx, t)$.

Having stated the weak divergence-free condition, we are prepared to prove \Cref{prop:div_free}, which states that among diffusion processes having a shared diffusion coefficient $\mbSigma^{1/2}$, the only degree of freedom is the choice of divergence-free vector field.

\begin{proof}[Proof of \Cref{prop:div_free}]
Let $\bbQ,\widetilde{\bbQ}\in\cQ$ have drifts $\mbf$ and
$\widetilde{\mbf}$, respectively, common diffusion coefficient
$\mbSigma^{1/2}$, and common initial marginal density $q(\mbx,0)$.
Let $\{q(\mbx,t)\}_{0\leq t\leq T}$ denote the marginal densities of
$\bbQ$.

Suppose first that $\bbQ$ and $\widetilde{\bbQ}$ have the same one-time marginals. 
Since these common marginals solve the weak Fokker--Planck equation for both drifts, and the diffusion coefficients are equal, subtracting the two equations gives
\begin{equation}\label{app:eqn:div-free-weak-diff}
    \int \left(\mbf(\mbx,t)-\widetilde{\mbf}(\mbx,t)\right) \cdot\nabla\varphi(\mbx) q(d\mbx,t) =0, \qquad\text{a.e.~$t$},
\end{equation}
for every test function $\varphi$. This is precisely the statement that $\mbf-\widetilde{\mbf}$ is $q$-divergence-free.

Conversely, suppose that $\mbf-\widetilde{\mbf}$ is $q$-divergence-free. 
Since $q$ solves the weak Fokker--Planck equation associated with $\mbf$, \cref{app:eqn:div-free-weak-diff} implies that it also solves the weak Fokker--Planck equation associated with $\widetilde{\mbf}$ and the same initial density. 
By \Cref{app:ass:fpk-unique}, this solution is unique. 
It therefore coincides with the marginal densities of $\widetilde{\bbQ}$, so $\bbQ$ and $\widetilde{\bbQ}$ have the same one-time marginals.
\end{proof}

There are also diffusion processes with different diffusion coefficients that share the same one-time marginals. This fact is leveraged by \citet{bartosh2025sde} to learn the shared prior and posterior diffusion coefficient. 
However, for a fixed prior SDE \cref{eqn:prior_sde}, the diffusion coefficient of the posterior \cref{eqn:variational_family} must be equal to that of the prior. 
This is because, by Girsanov's theorem (\Cref{app:thm:girsanov}), diffusion processes whose diffusion coefficients differ are generally not mutually absolutely continuous. 
In other words, the diffusion coefficient must be viewed as a parameter of the generative model, rather than a degree of freedom in the posterior process. 
On the other hand, any two diffusion processes with shared diffusion coefficient will be mutually absolutely continuous.

\subsection{Gaussian one-time marginals}
\label{app:subsec:gaussian}

In the case of Gaussian one-time marginals, the non-uniqueness of the path law described in \Cref{app:subsec:div-free} can be made explicit. In this subsection, we provide a characterization of the affine vector fields that are $q$-divergence free.

Let $(\mbm(t),\mbS(t))_{0\leq t\leq T}$ be continuously differentiable and $\mbS(t)$ be positive definite. 
For a matrix-valued function $\mbK(t)$, define the affine velocity field
\begin{equation}\label{eq:gaussian-affine-velocity}
    \mbv_{\mbK}(\mbx,t) = \frac{d}{d t}\mbm(t) + \mbK(t)(\mbx- \mbm(t)).
\end{equation}
And consider the SDE having diffusion coefficient $\mbSigma(\mbx,t)^{1/2}$ and drift
\begin{equation}\label{app:eq:gaussian-compatible-drift}
\begin{aligned}
    &\mbf_{q}(\mbx,t) = \mbv_{\mbK}(\mbx,t) + \frac{1}{2}\mbSigma(\mbx,t)\nabla_{\mbx}\log q(\mbx,t) + \frac{1}{2}\sum_j \partial_{x_j}\mbSigma_{\cdot,j}(\mbx,t)\\
    &= \frac{d}{dt}\mbm(t) + \mbK(t)(\mbx-\mbm(t)) - \frac{1}{2}\mbSigma(\mbx,t)\mbS(t)^{-1}(\mbx-\mbm(t)) + \frac{1}{2}\sum_j \partial_{x_j}\mbSigma_{\cdot,j}(\mbx,t).
\end{aligned}
\end{equation}

In \Cref{app:prop:gaussian-compatible-drift}, we characterize the choices of $\mbK(t)$ for which this SDE produces one-time marginals $q(\mbx,t) = \cN(\mbm(t),\mbS(t))$. 
\begin{proposition}
\label{app:prop:gaussian-compatible-drift}
Fix any continuously differentiable square-root factorization $\mbL(t)$ satisfying
$\mbL(t)\mbL(t)^\top=\mbS(t)$. Then the following statements are equivalent:
\begin{enumerate}[label=(\roman*)]
    \item The SDE
    \begin{equation*}
        d\mbx(t) = \mbf_{q}(\mbx(t),t)dt + \mbSigma(\mbx(t),t)^{1/2}d\mbw(t), \quad q(\mbx, 0) =  \cN(\mbm(0),\mbS(0))
    \end{equation*}
    has one-time marginals $q(\mbx,t)=\cN(\mbm(t),\mbS(t))$.

    \item $\mbK(t)$ satisfies the matrix Lyapunov equation
    \begin{equation} \label{app:eq:gaussian-lyapunov}
        \frac{d}{d t}\mbS(t) = \mbK(t)\mbS(t) + \mbS(t)\mbK(t)^\top, \qquad 0\leq t\leq T.
    \end{equation}

    \item There exists a unique skew-symmetric matrix-valued function
    $\mbOmega(t)=-\mbOmega(t)^\top$ such that
    \begin{equation}\label{app:eq:gaussian-square-root}
        \mbK(t) = \left(\frac{d}{d t}\mbL(t)\right)\mbL(t)^{-1} + \mbL(t)\mbOmega(t)\mbL(t)^{-1}, \qquad 0\leq t\leq T.
    \end{equation}
\end{enumerate}
\end{proposition}

\begin{proof}
Before proving the equivalence of (i), (ii), and (iii), we show that the Fokker--Planck equation associated with
the drift \eqref{app:eq:gaussian-compatible-drift} and diffusion coefficient $\mbSigma(\mbx, t)^{1/2}$ reduces to a continuity equation. 
Using $q(\mbx, t)\nabla_{\mbx}\log q(\mbx, t)=\nabla_{\mbx} q(\mbx, t)$,
\begin{equation*}
\begin{aligned}
    q(\mbx,t)\mbf_{q}(\mbx,t) &= q(\mbx,t)\mbv_{\mbK}(\mbx,t) + \frac{1}{2}\mbSigma(\mbx,t)\nabla_{\mbx}q(\mbx,t) + \frac{1}{2}q(\mbx,t) \sum_j\partial_{x_j}\mbSigma_{\cdot,j}(\mbx,t)\\
    &= q(\mbx,t)\mbv_{\mbK}(\mbx,t) + \frac{1}{2}\sum_j \partial_{x_j}\left[ q(\mbx,t)\mbSigma_{\cdot,j}(\mbx,t)\right].
\end{aligned}
\end{equation*}
Consequently, we can rewrite the Fokker--Planck equation as
\begin{align}
    \frac{\partial}{\partial t}q(\mbx,t) &= -\nabla_{\mbx}\cdot \left[ q(\mbx,t)\mbf_{q}(\mbx,t) \right] + \frac{1}{2} \sum_{i,j} \partial_{x_i,x_j} \left[q(\mbx,t)\mbSigma_{ij}(\mbx,t)\right]\nonumber\\
    &= -\nabla_{\mbx}\cdot\left[q(\mbx,t)\mbv_{\mbK}(\mbx,t) \right].\label{app:eq:continuity-affine}
\end{align}
This is the continuity equation corresponding to the affine flow
\begin{equation}\label{app:eq:affine-flow}
    \frac{d}{d t}\mbx(t)
    =
    \mbv_{\mbK}(\mbx(t),t).
\end{equation}
Because the flow map is affine, it maps the Gaussian initial distribution $q(\mbx,0)=\cN(\mbm(0),\mbS(0))$ to a Gaussian distribution at every time.

To prove the equivalence of (i) and (ii), we must show that  $q(\mbx,t)=\cN(\mbm(t),\mbS(t))$ solves \cref{app:eq:continuity-affine} if and only if $\mbK(t)$ satisfies the Lyapunov equation \cref{app:eq:gaussian-lyapunov}.

Let $\widetilde{\mbm}(t) = \bbE[\mbx(t)]$
denote the mean induced by the affine flow \cref{app:eq:affine-flow}. 
By the weak form of the continuity equation \cref{app:eq:continuity-affine}, $\widetilde{\mbm}(t)$ satisfies the ODE
\begin{equation*}
    \frac{d}{d t}\widetilde{\mbm}(t) = \frac{d}{d t}\mbm(t) + \mbK(t) \left(\widetilde{\mbm}(t)-\mbm(t) \right), \quad 
    \widetilde{\mbm}(0) = \mbm(0).
\end{equation*}
It is immediate that $\mbm(t)$ satisfies the ODE for the mean of $\mbx(t)$. 
For the covariance, we define the centered state $\mbz(t) = \mbx(t)-\mbm(t)$, which satisfies the affine flow
\begin{equation*}
    \frac{d}{d t}\mbz(t) = \mbK(t)\mbz(t).
\end{equation*}
Therefore, if $\widetilde{\mbS}(t) = \bbE \left[ \mbz(t)\mbz(t)^\top \right]$ denotes the covariance induced by the affine flow, the weak form of continuity equation yields the ODE
\begin{equation*}
    \frac{d}{d t}\widetilde{\mbS}(t) =\bbE\left[ \mbK(t)\mbz(t)\mbz(t)^\top + \mbz(t)\mbz(t)^\top\mbK(t)^\top\right] = \mbK(t)\widetilde{\mbS}(t) + \widetilde{\mbS}(t)\mbK(t)^\top, \quad  \widetilde{\mbS}(0)=\mbS(0).
\end{equation*}
If $\mbK(t)$ satisfies the Lyapunov equation
\cref{app:eq:gaussian-lyapunov}, uniqueness of solutions to the covariance ODE implies $\widetilde{\mbS}(t)=\mbS(t)$. 

Conversely, if $q(\mbx,t)=\cN(\mbm(t),\mbS(t))$ solves
\cref{app:eq:continuity-affine}, then its covariance must satisfy the same equation, so
\begin{equation*}
    \frac{d}{d t}\mbS(t) = \mbK(t)\mbS(t) + \mbS(t)\mbK(t)^\top.
\end{equation*}
This proves the equivalence of \textnormal{(i)} and \textnormal{(ii)}.

Suppose next that $\mbK(t)$ has the representation in
\eqref{app:eq:gaussian-square-root}. Since
$\mbS(t)=\mbL(t)\mbL(t)^\top$, we obtain
\begin{equation*}
\begin{aligned}
    \mbK(t)\mbS(t)+\mbS(t)\mbK(t)^\top &= \left(\frac{d}{d t}\mbL(t) \right)\mbL(t)^\top + \mbL(t) \left(\frac{d}{t}\mbL(t)  \right)^\top + \mbL(t) \left( \mbOmega(t)+\mbOmega(t)^\top\right)\mbL(t)^\top\\
    &= \frac{d}{d t} \left[\mbL(t)\mbL(t)^\top\right] =
    \frac{d}{d t}\mbS(t).
\end{aligned}
\end{equation*}
Thus, (iii) implies (ii).

Conversely, suppose that $\mbK(t)$ satisfies \eqref{app:eq:gaussian-lyapunov} and define
\begin{equation*}
    \mbOmega(t) = \mbL(t)^{-1}\mbK(t)\mbL(t) - \mbL(t)^{-1}\frac{d}{d t}\mbL(t).
\end{equation*}
Then
\begin{equation*}
    \mbOmega(t) + \mbOmega(t)^\top  = \mbL(t)^{-1} \left[    \mbK(t)\mbS(t) + \mbS(t)\mbK(t)^\top - \frac{d}{d t}\mbS(t) \right] \mbL(t)^{-\top} =\mb{0}.
\end{equation*}
Hence, $\mbOmega(t)$ is skew-symmetric. 
Its definition also shows that it is unique for the fixed square root $\mbL(t)$, proving the equivalence of (ii) and (iii).
\end{proof}

For a fixed square root $\mbL(t)$, the skew-symmetric matrix $\mbOmega(t)$ parameterizes the
freedom among compatible affine velocities. In particular, changing $\mbOmega(t)$ adds the vector
field
\begin{equation*}
    \mbh_{\mbOmega}(\mbx,t) =  \mbL(t)\mbOmega(t)\mbL(t)^{-1}(\mbx - \mbm(t)),
\end{equation*}
which satisfies
\begin{equation*}
    \nabla_{\mbx}\cdot \left[ q(\mbx,t)\mbh_{\mbOmega}(\mbx,t) \right] = 0.
\end{equation*}
Thus, $\mbh_{\mbOmega}$ changes the path law without changing the one-time marginals. 
When $K\geq2$, the space of skew-symmetric matrices has dimension $\binom{K}{2}$.

The representation in \eqref{app:eq:gaussian-square-root} depends on the chosen square root $\mbL(t)$.
Replacing $\mbL(t)$ by another factor of the covariance $\mbS(t)$ reparameterizes $\mbOmega(t)$ but does not change the set of compatible $\mbK(t)$.

Finally, Proposition~\ref{app:prop:gaussian-compatible-drift} characterizes only the subfamily of affine $q$-divergence-free vector fields. 
There also exist nonlinear $q$-divergence-free vector fields that preserve Gaussian one-time marginals while inducing non-Gaussian multi-time distributions.

% \subsection{Gaussian one-time marginals}\label{app:subsec:gaussian}
% In the case of Gaussian one-time marginals, the non-uniqueness of the path law described in \Cref{app:subsec:div-free} can be made explicit. 
% In this subsection, we construct an infinite collection of (distinct) path laws whose one-time marginals are identical. Throughout, we assume $q(\mbx, t) = \cN(\mbm(t), \mbS(t))$, where $\mbm(t), \mbS(t)$ are continuously differentiable on $[0, T]$ and $\mbS$ is positive definite.

% We first characterize these path laws in terms of the solution to a matrix Lyapunov equation. 
% \begin{proposition}\label{prop:lyapunov}
% Let $q(\mbx, t) = \cN(\mbm(t), \mbS(t)), 0 \leq t \leq T$ be a family of Gaussian one-time marginal distributions, with $(\mbm(t), \mbS(t))_{0 \leq t \leq T}$ continuously differentiable in $t$ and $\mbS(t)$ positive definite. Let $(\mbK(t))_{0 \leq t \leq T}$ be any solution to the matrix Lyapunov equation
% \begin{equation}\label{app:eqn:lyapunov}
%     \frac{d}{dt} \mbS(t) = \mbK(t) \mbS(t) + \mbS(t) \mbK(t)^\top, \quad 0 \leq t \leq T.
% \end{equation}
% Then the SDE with drift $\mbf_q(\mbx, t) = \mbA_q(\mbx, t) \mbx + \mbb_q(\mbx, t) + \frac{1}{2}\sum_j \partial_{\mbx_j} \mbSigma_{\cdot, j}(\mbx, t)$,
% \begin{equation}\label{app:eqn:lyapunov-drift}
%     \mbA_q(\mbx, t) = \mbK(t) - \frac{1}{2}\mbSigma(\mbx, t) \mbS(t)^{-1}, \quad \mbb_q(\mbx, t) = \frac{d}{dt}\mbm(t) - \mbA_q(\mbx, t)\mbm(t),
% \end{equation}
% and diffusion coefficient $\mbSigma(\mbx, t)^{1/2}$ has marginals $q(\mbx, t) = \cN(\mbm(t), \mbS(t))$, $0 \leq t \leq T$.
% \end{proposition}
% The Lyapunov equation \eqref{app:eqn:lyapunov} has infinitely many solutions $\mbK(t)$, so a single collection of  Gaussian one-time marginals $\{ q(\mbx, t) \}$ corresponds to infinitely many compatible diffusions. 
% Note that \Cref{prop:lyapunov} does \emph{not} characterize the complete set of diffusions consistent with Gaussian one-time marginals. 
% For example, when $\mbSigma$ is state-independent, \Cref{prop:lyapunov} only describes Gaussian processes i.e., stochastic process whose multi-time joint distributions are multivariate Gaussian. There exist stochastic processes with Gaussian one-time marginals that are not Gaussian processes.

% The explicit construction is obtained by plugging the drift from \Cref{prop:lyapunov} into the Fokker-Planck equation, which yields the \emph{continuity equation} \cref{app:eqn:fpk-simple} for Gaussian marginals distributions and a linear vector field. 
% Unlike the Fokker-Planck equation, the continuity equation does not depend on the diffusion coefficient $\mbSigma^{1/2}$, which is a consequence of the two terms in the posterior drift cancelling with the diffusion term in the Fokker-Planck equation. 
% The same strategy is used to convert between ODE and SDE samplers for diffusion and flow matching models \cite{bartosh2025sde,songscore2021, albergo2025stochastic}.See also \Cref{app:subsec:svise-sde-compare}.

% \begin{proof}[Proof of \Cref{prop:lyapunov}]
% Note that since $q(\mbx, t)$ and $\mbf_q$ are sufficiently differentiable, we can work directly with the strong form of the Fokker--Planck equation \cref{eqn:fpk}.

% Starting first with the diffusion term, we have by the chain rule and the formula for the Gaussian density
% {\small
% \begin{equation*}
%     \begin{aligned}
%     \frac{1}{2}\sum_{i, j} \partial_{\mbx_i, \mbx_j} [q(\mbx, t) \mbSigma_{ij}(\mbx,t)]
%     &= \frac{1}{2}\sum_i \partial_{\mbx_i}
%     \left[
%     q(\mbx,t)\sum_j
%     \left\{
%     \partial_{\mbx_j} \mbSigma_{ij}(\mbx,t)
%     -
%     [\mbS(t)^{-1}(\mbx-\mbm(t))]_j \mbSigma_{ij}(\mbx,t)
%     \right\}
%     \right] \\
%     &= \frac{1}{2}\nabla_{\mbx} \cdot \left[
%     - q(\mbx, t)\mbSigma(\mbx,t) \mbS(t)^{-1}(\mbx-\mbm(t))
%     + q(\mbx, t) \sum_j \partial_{\mbx_j} \mbSigma_{\cdot, j}(\mbx,t)
%     \right].
%     \end{aligned}
% \end{equation*}
% }
% Notice that the same terms with the opposite sign appear in the drift term of the Fokker--Planck equation, by our definition of the drift $\mbf_q$. Indeed,
% {\small
% \begin{equation*}
% \begin{aligned}
%   &-\nabla_{\mbx} \cdot [\mbf_q(\mbx, t) q(\mbx, t)]
%   =\\
%   &-\nabla_{\mbx} \cdot \bigg[
%    \bigg(
%    \mbK(t)(\mbx-\mbm(t))
%    + \frac{d}{dt}\mbm(t)
%    - \frac{1}{2}\mbSigma(\mbx,t)\mbS(t)^{-1}(\mbx-\mbm(t))
%    + \frac{1}{2} \sum_j \partial_{\mbx_j}\mbSigma_{\cdot,j}(\mbx,t)
%    \bigg) q(\mbx,t)
%    \bigg],
% \end{aligned}
% \end{equation*}
% }
% where we used
% \begin{equation*}
%     \mbA_q(\mbx,t)\mbx+\mbb_q(\mbx,t)
%     =
%     \mbA_q(\mbx,t)(\mbx-\mbm(t))+\frac{d}{dt}\mbm(t).
% \end{equation*}
% Canceling these terms, this shows that the Fokker--Planck equation reduces to the continuity equation
% \begin{equation}\label{app:eqn:fpk-simple}
%     \frac{\partial q(\mbx, t)}{\partial t}
%     =
%     -\nabla_{\mbx} \cdot
%     \left[
%     \left(
%     \mbK(t)(\mbx-\mbm(t))
%     + \frac{d}{dt}\mbm(t)
%     \right) q(\mbx, t)
%     \right],
%     \quad 0 \leq t \leq T.
% \end{equation}

% Next, computing both sides of \Cref{app:eqn:fpk-simple} explicitly and dividing by $q(\mbx,t)$, we obtain
% {\small
% \begin{equation*}
% \begin{aligned}
% &-\frac{1}{2}\bigg[
% \mathrm{tr}\left(\mbS(t)^{-1} \frac{d}{dt} \mbS(t)\right)
% - (\mbx - \mbm(t))^\top  \mbS(t)^{-1} \frac{d}{dt} \mbS(t) \mbS(t)^{-1} (\mbx - \mbm(t))
% - 2 \frac{d}{dt}\mbm(t)^\top  \mbS(t)^{-1}  (\mbx - \mbm(t))
% \bigg]\\
% &=  -\left[
% \mathrm{tr}(\mbK(t))
% -
% \left(
% \mbK(t)(\mbx-\mbm(t))
% + \frac{d}{dt}\mbm(t)
% \right)^\top
% \mbS(t)^{-1}(\mbx - \mbm(t))
% \right].
% \end{aligned}
% \end{equation*}
% }
% The linear terms involving $\frac{d}{dt}\mbm(t)$ cancel, so this further simplifies to
% \begin{equation}\label{app:eqn:fpk-equality}
% \begin{aligned}
% &-\frac{1}{2}\left[
% \mathrm{tr}\left(\mbS(t)^{-1} \frac{d}{dt} \mbS(t)\right)
% - (\mbx - \mbm(t))^\top  \mbS(t)^{-1} \frac{d}{dt} \mbS(t) \mbS(t)^{-1}(\mbx - \mbm(t))
% \right]\\
% =& -\left[
% \mathrm{tr}(\mbK(t))
% - (\mbx - \mbm(t))^\top \mbK(t)^\top \mbS(t)^{-1}(\mbx - \mbm(t))
% \right].
% \end{aligned}
% \end{equation}

% Now, we must determine conditions on $(\mbK(t))_{0 \leq t \leq T}$ under which equality \eqref{app:eqn:fpk-equality} holds. If $\mbK(t)$ solves the Lyapunov equation \eqref{app:eqn:lyapunov}, then
% \begin{equation*}
% \begin{aligned}
%     &\mbS(t)^{-1} \frac{d}{dt} \mbS(t) \mbS(t)^{-1}
%     =
%     \mbS(t)^{-1}\mbK(t) + \mbK(t)^\top\mbS(t)^{-1} \\
%     \implies&
%     (\mbx - \mbm(t))^\top
%     \mbS(t)^{-1} \frac{d}{dt} \mbS(t) \mbS(t)^{-1}
%     (\mbx - \mbm(t)) =
%     2(\mbx - \mbm(t))^\top
%     \mbK(t)^\top\mbS(t)^{-1}
%     (\mbx - \mbm(t)).
% \end{aligned}
% \end{equation*}
% Moreover, $\mathrm{tr}(\mbS(t)^{-1} \frac{d}{dt}\mbS(t)) = 2\mathrm{tr}(\mbK(t))$. Hence, whenever $\mbK(t)$ solves the Lyapunov equation, the diffusion process with drift $\mbf_q$ and diffusion coefficient $\mbSigma^{1/2}$ satisfies the Fokker--Planck equation.
% \end{proof}


% To conclude this subsection, we provide an alternate (but equivalent) characterization of the Lyapunov equation \cref{app:eqn:lyapunov} for SDEs that induce the same one-time Gaussian marginals.
% We do so via square root factorizations of $\mbS(t)$.
% \begin{proposition}\label{app:prop:lyapunov-gauge}
% Fix any continuously differentiable square root factorization $\mbL(t)$ satisfying $\mbL(t)\mbL(t)^\top = \mbS(t)$. Then $\mbK(t)$ solves the Lyapunov equation \eqref{app:eqn:lyapunov} if and only if
% \begin{equation}\label{app:eqn:lyapunov-gauge}
% \mbK(t) = \left(\frac{d}{dt}\mbL(t) \right)\mbL(t)^{-1} + \mbL(t)\mbOmega(t)\mbL(t)^{-1}, \quad 0 \leq t \leq T
% \end{equation}
% for some skew-symmetric matrix-valued function $\mbOmega(t) = -\mbOmega(t)^\top$.
% \end{proposition}

% \begin{proof}
% \textbf{Sufficiency.}
% Suppose is $\mbK$ defined as in  \cref{app:eqn:lyapunov-gauge}. Substituting this expression into the right-hand side of the Lyapunov equation \cref{app:eqn:lyapunov} yields
% \begin{equation*}
% \begin{aligned}
% \mbK\mbS + \mbS\mbK^\top 
% &= \left( \left(\frac{d}{dt}\mbL\right)\mbL^{-1} + \mbL\mbOmega\mbL^{-1}\right)\mbL\mbL^\top + \mbL\mbL^\top\left(\mbL^{-\top}\left(\frac{d}{dt}\mbL\right)^\top + \mbL^{-\top}\mbOmega^\top \mbL^\top \right)\\
% &= \left(\frac{d}{dt}\mbL\right)\mbL^\top + \mbL\mbOmega\mbL^\top + \mbL\left(\frac{d}{dt}\mbL\right)^\top + \mbL\mbOmega^\top\mbL^\top \\
% &\overset{(\star)}{=} \left(\frac{d}{dt}\mbL\right)\mbL^\top + \mbL\left(\frac{d}{dt}\mbL\right)^\top\\
% &= \frac{d}{dt}(\mbL\mbL^\top) = \frac{d}{dt}\mbS.
% \end{aligned}
% \end{equation*}
% Equality $(\star)$ holds by skew-symmetry $\mbOmega + \mbOmega^\top = \mb{0}$.

% \textbf{Necessity.} Now suppose $\mbK$ solves the Lyapunov equation. Solving \cref{app:eqn:lyapunov-gauge} for $\mbOmega$ yields  $\mbOmega = \mbL^{-1}\left(\mbK - \left(\frac{d}{dt} \mbL \right) \mbL^{-1} \right)\mbL$. 
% Hence, we need to show $\mbOmega$ is skew-symmetric.
% A direct calculation reveals
% \begin{equation*}
% \begin{aligned}
% \mbOmega + \mbOmega^\top &= \mbL^{-1}\left(\mbK\mbS - 
% \left(\frac{d}{dt}\mbL \right)\mbL^\top \right)\mbL^{-\top} + \mbL^{-1}\left(\mbS\mbK^\top - \mbL\left(\frac{d}{dt}\mbL \right)^\top\right)\mbL^{-\top}\\
% &= \mbL^{-1}\left(\mbK\mbS + \mbS\mbK^\top - \frac{d}{dt} \mbS \right)\mbL^{-\top} = \mb{0}.
% \end{aligned}
% \end{equation*}
% \end{proof}

% The two terms in \eqref{app:eqn:lyapunov-gauge} have distinct roles. 
% $(\frac{d}{dt} \mbL)\mbL^{-1}$ depends on the chosen square root and is itself a solution to the Lyapunov equation. 
% Arbitrarily changing $\mbL$ will change $\mbL \mbL^\top$ and hence the one-time marginals.
% On the other hand, $\mbL\mbOmega\mbL^{-1}$ is allowed to vary as $\mbOmega$ takes different values in the $\binom{K}{2}$-dimensional vector space of skew-symmetric matrices. For fixed $\mbL$, each choice of $\mbOmega$ yields the same one-time marginals.

\subsection{SVISE (\citet{course2023state}) and SDE Matching (\citet{bartosh2025sde})}\label{app:subsec:svise-sde-compare}

In this subsection, we explain how existing simulation-free VI methods SVISE \cite{course2023state, course2023amortized} and SDE Matching \cite{bartosh2025sde} select a single drift function that realizes a given collection of one-time Gaussian marginal distributions $q(\mbx, t) = \cN(\mbm(t), \mbS(t))$, where $\mbm(t), \mbS(t)$ are continuously differentiable on $[0, T]$ and $\mbS(t)$ is positive definite. 
We refer to these choices as the symmetric and square root gauges, respectively.

Before describing each choice of posterior drift, we recall \Cref{app:prop:gaussian-compatible-drift} from \Cref{app:subsec:gaussian}, which tells us that
\begin{equation}\label{app:eqn:variational-drift}
\mbf_q(\mbx, t) = \frac{d}{dt}\mbm(t) + \mbK(t)(\mbx-\mbm(t)) - \frac{1}{2}\mbSigma(\mbx,t)\mbS(t)^{-1}(\mbx-\mbm(t)) + \frac{1}{2}\sum_j \partial_{x_j}\mbSigma_{\cdot,j}(\mbx,t)
\end{equation}
attains the marginals $q(\mbx, t)$ whenever $\mbK(t)$ is a solution to the Lyapunov equation \cref{app:eq:gaussian-lyapunov}.
As we will see, SVISE and SDE Matching select the variational drift \cref{app:eqn:variational-drift} using different choices of $\mbK$.

\paragraph{Symmetric gauge.}
Suppose $\mbSigma(t) = \mbSigma(\mbx, t)$ is state-independent and $\mbf_q$ is defined as in \Cref{app:prop:gaussian-compatible-drift}. 
Then  $\mbA_q(t) = \mbK(t) - \frac{1}{2}\mbSigma(t)\mbS(t)^{-1}$ also satisfies a Lyapunov equation
\begin{equation}\label{eqn:course-nair}
    \frac{d}{dt} \mbS(t) - \mbSigma(t)  = \mbA_q(t)\mbS(t) + \mbS(t) \mbA_q(t)^\top, \quad 0 \leq t \leq T.
\end{equation}

In SVISE \cite{course2023state, course2023amortized}, the authors choose the posterior drift $\mbA_q(t)$ to be the unique symmetric solution to \cref{eqn:course-nair}. They write this solution in closed-form using the Kronecker sum.

Unlike the square root gauge, the symmetric gauge does not naturally extend to the case of a state-dependent diffusion coefficient. This is because, in general, the matrix $\mbK$ for which $\mbA_q(\mbx, t) = \mbK - \frac{1}{2}\mbSigma(\mbx, t) \mbS(t)^{-1}$ is symmetric cannot be chosen to be independent of the state $\mbx$.

\paragraph{Square root gauge.}
\citet{bartosh2025sde} instead identify the one-time Gaussian marginals with the drift function $\mbf_q$ defined as in \Cref{app:prop:gaussian-compatible-drift} for $\mbK(t) = \bigl(\frac{d}{dt}\mbS(t)^{1/2}\bigr)\mbS(t)^{-1/2}$. 
Here, $\mbS(t)^{1/2}$ denotes the symmetric square root of $\mbS(t)$.
It is straightforward to verify that this choice of $\mbK$ satisfies the matrix Lyapunov equation \eqref{app:eq:gaussian-lyapunov}. This choice of $\mbK$ is valid even for a state-dependent diffusion coefficient.

The SDE Matching framework, in theory, allows for any square root $\mbsigma(t)$ such that $\mbS(t) = \mbsigma(t) \mbsigma(t)^\top$, to be chosen. 
However, in their implementation, $\mbsigma(t)$ is chosen to be diagonal, which yields the symmetric square root $\mbS(t)^{1/2}$.
At the end of this section, we discuss an extension of SDE Matching that represents all square roots simultaneously.

The authors arrive at the choice of variational drift by converting from a probability flow ODE \cite{songscore2021} to an SDE. In particular, suppose $\mbvarepsilon \sim \cN(\mb{0}, \bbI_K)$ and $F(\mbvarepsilon, t) = \mbm(t) + \mbS(t)^{1/2}\mbvarepsilon$.
$F$ can be viewed as a map that takes as input a sample $\mbvarepsilon$ as well as a time $t$ and produces a sample from $q(\mbx, t)$.
If $\tilde{\mbx}(t) = F(\mbvarepsilon, t)$, then each particle $\tilde{\mbx}(t)$ evolves according to the ODE
\begin{equation}\label{app:eqn:pf-ode}
    \begin{aligned}
    \frac{d}{dt} \tilde{\mbx}(t) &= \frac{d}{dt}\mbm(t) + \left(\frac{d}{dt}\mbS(t)^{1/2}\right)\mbvarepsilon\\
    &= \frac{d}{dt}\mbm(t) + \left(\frac{d}{dt}\mbS(t)^{1/2}\right)\{\mbS(t)^{-1/2}(\tilde{\mbx}(t) - \mbm(t))\}, \quad 0 \leq t \leq T
    \end{aligned}
\end{equation}
with $\tilde{\mbx}(0) \overset{d}{=} F(\mbvarepsilon, 0)$. 

From the ODE \eqref{app:eqn:pf-ode}, the authors construct an SDE with diffusion coefficient $\mbSigma(\mbx, t)^{1/2}$. A straightforward manipulation of the Fokker-Planck equation (see e.g., \citet[][Appendix D.1]{songscore2021}) shows that the solution to the SDE
\begin{equation*}
\begin{aligned}
    d \mbx(t) &= \bigg( \frac{d}{dt}\mbm(t) + \left(\frac{d}{dt}\mbS(t)^{1/2}\right)\{\mbS(t)^{-1/2}(\mbx(t) - \mbm(t))\} + \frac{1}{2} \mbSigma(\mbx, t) \nabla_{\mbx} \log q(\mbx(t), t)\\
    &+ \frac{1}{2} \sum_{j} \partial_{\mbx_j} \mbSigma_{\cdot, j}(\mbx(t), t) \bigg)dt + \mbSigma(\mbx(t), t)^{1/2} d\mbw(t), \quad \mbx(0) \overset{d}{=} q(\mbx, 0)
\end{aligned}
\end{equation*}
has the same one-time marginal distributions as $\tilde{\mbx}(t)$. Substituting in the Gaussian score and grouping terms yields the drift
\begin{equation*}
\begin{aligned}
   &\left\{ \left(\frac{d}{dt}\mbS(t)^{1/2}\right)\mbS(t)^{-1/2} - \frac{1}{2}\mbSigma(\mbx(t), t)\mbS(t)^{-1}\right\}\mbx(t) +\\
    &\frac{d}{dt}\mbm(t) - \left\{ \left(\frac{d}{dt}\mbS(t)^{1/2}\right)\mbS(t)^{-1/2} - \frac{1}{2}\mbSigma(\mbx(t), t)\mbS(t)^{-1}\right\} \mbm(t) + \frac{1}{2} \sum_{j} \partial_{\mbx_j} \mbSigma_{\cdot, j}(\mbx(t), t).
\end{aligned}
\end{equation*}
Once we identify $\mbK(t) = \left(\frac{d}{dt}\mbS(t)^{1/2}\right)\mbS(t)^{-1/2}$, the resulting drift is of the form given in \cref{app:eq:gaussian-compatible-drift}.

Although the construction in \citet{bartosh2025sde} allows the marginals $q(\mbx, t)$ to be non-Gaussian (e.g., by parameterizing $F$ as a normalizing flow), the paper itself evaluates the method only in the Gaussian marginal setting. We therefore restrict our discussion to that case.


\paragraph{An extension of SDE Matching.}
The framework of SDE Matching \cite{bartosh2025sde} can in principle be extended to span every solution to the Lyapunov equation \eqref{app:eq:gaussian-lyapunov}.
This is not discussed or implemented by the authors. 
As described in the previous paragraph, their parameterization $\mbS(t) = \mbsigma(t)\mbsigma(t)^\top$ induces the linear coefficient $\mbK(t) = (\frac{d}{dt}\mbsigma(t)) \mbsigma(t)^{-1}$. 
Fixing $\mbsigma$ to a specific square root selects a single path law for each set of one-time marginals. 
By \Cref{app:prop:gaussian-compatible-drift}, the full set of solutions to the Lyapunov equation is parameterized by an additional skew-symmetric matrix $\mbOmega(t)$ via
\begin{equation*}
    \mbK(t) = \left(\frac{d}{dt}\mbsigma(t) \right) \mbsigma(t)^{-1} + \mbsigma(t) \mbOmega(t) \mbsigma(t)^{-1}, \quad 0 \leq t \leq T.
\end{equation*}
Augmenting SDE Matching with a learned $\mbOmega(t)$ would in principle allow the posterior to span the full $\binom{K}{2}$-dimensional space of compatible linear drifts.

This extension has two limitations. 
First, it requires introducing and optimizing over $\mbOmega(t)$, sacrificing the closed-form structure of SDE Matching's original parameterization. 
Helmholtz-SDE, by contrast, \emph{computes the optimal correction in closed form} via the Helmholtz decomposition of the residual vector field (\Cref{sec:method}). 
Second, for a state-independent diffusion coefficient, the resulting path law remains a Gaussian process.
\Cref{app:prop:gaussian-compatible-drift} characterizes only diffusions with affine drift that produce the prescribed Gaussian marginals; non-linear drifts can also yield Gaussian one-time marginals while inducing non-Gaussian multi-time joint distributions. 
Helmholtz-SDE can produce posterior approximations that lie outside the Gaussian-process family by using higher-order polynomial approximations to the optimal drift (\Cref{app:subsec:resid-approx}).